\section{Foundations}\label{part:foundations}

This section assembles the search core from its parts: the data model every other
piece reads and writes, the propagation that reasons over it, and the branching
that turns a pruned domain into a concrete solution. We start with how a variable
is stored, look inside a single propagator, sort out which high-level constraints
are expanded away versus kept as propagators, and finish with how search drives a
range down to one value.

\begin{platypusbook}[htb]
A working familiarity with CDCL SAT solving (decisions, unit propagation,
conflict analysis, clause learning, and backjumping) is not strictly required
here, but it is highly recommended. CP-SAT layers a good deal on top of a CDCL
core, so having seen how a plain CDCL solver works makes the rest considerably
easier to follow. Readers who want that background first can start with Smock's
gentle visual introduction
\cite{smock2016peek} (about 35 minutes), move on to Nadel's more technical talk
linking CDCL to optimization \cite{nadel2023cdcl}, and turn to Biere's in-depth
tutorial \cite{biere2021satsolving} (about four and a half hours) for the full
machinery. Knuth's fascicle \cite{knuth2015satisfiability} is the canonical
written reference, useful either in place of the videos or alongside them.
\end{platypusbook}

\subsection{The data model: variables as bounds}\label{sec:bounds}
% Sources (VERIFIED):
%   ortools/sat/integer_base.h -- IntegerLiteral (x >= bound / x <= bound),
%     GreaterOrEqual()/LowerOrEqual(); negation flips the bound.
%   ortools/sat/integer.h -- IntegerTrail stores monotone lower bounds
%     (upper bound = lower bound on NegationOf(var)); a var is fixed when lb==ub.

Everything that follows rests on a single design decision: \emph{how the solver
represents an integer variable while it is running}. Once this idea is clear,
the rest of the architecture (propagation, reasons, and learning) falls into
place almost inevitably.

\textbf{Variables are bounds, not values.} An integer variable $x$ with domain
$[0,10]$ is \emph{not} stored as a tentative value, such as ``$x$ is currently
$3$.'' It carries no working assignment at all: nothing like the fractional
value a MIP branch-and-bound node rounds, or the running guess a heuristic nudges
toward feasibility. Instead it is stored as a pair of \textbf{current bounds}, a
lower bound $\text{lb}(x)$ and an upper bound $\text{ub}(x)$. Search and
propagation never \emph{assign} $x$; they only \emph{narrow} the gap between these
two numbers, and $x$ is \emph{fixed} only once they meet.

This puts CP-SAT close to a classical finite-domain solver, which also narrows
rather than assigns, but with one decisive difference. A finite-domain solver
keeps the full set of surviving candidate values and lets a propagator strike any
one of them, every even number say, out of the interior of a domain. CP-SAT's
propagators move only the lower and upper bounds. The integer trail reasons about
$x$ through those two numbers alone, treating the live domain as the interval
$[\text{lb}(x), \text{ub}(x)]$ between them. When the declared domain has holes,
this interval is an \emph{over-approximation}: bound propagation works on the
enclosing interval, while the holes themselves are tracked separately, by presolve
canonicalization and the integer encoding of \cref{sec:encoding}. Until that
section, treat every domain as a plain interval.

As the solver descends the search tree, the bounds tighten ($\text{lb}$ rises and
$\text{ub}$ falls); when it backtracks, they relax again. Domains therefore shrink
from the outside in and grow back the same way.

Two properties of this representation make it work:

\begin{itemize}
  \item \textbf{Monotonicity.} Within one root-to-node path, the bounds only move
  inward. A propagator never has to reconcile competing updates; it only reads
  the tightest bounds known so far.
  \item \textbf{Reversibility.} Because a bound is a single number stamped with
  the decision level that set it, undoing work on backtrack is simple: pop every
  bound change made above the level to which the solver jumps. This integer
  trail, a stack of bound changes, is the same one that powers conflict analysis
  (\cref{sec:learning}).
\end{itemize}

\textbf{The atomic fact: a bound literal.} If variables are bounds, then the
elementary statements the solver asserts and reasons about are \textbf{bound
literals}: claims of the form
\[
  \bl{x \ge 5} \qquad\text{or}\qquad \bl{x \le 3} .
\]
The double brackets mark the claim as an \emph{object} the solver holds,
enqueues, and gathers into reasons, as opposed to the bare arithmetic relation
$x \ge 5$; we write them where that literal-as-object reading is what matters and
drop them in running prose where it is not. A bound literal is the integer analog
of a Boolean SAT literal, and, crucially, it behaves like one in the two ways the
learning machinery needs. (These are
exactly what the source calls an \srcla{integer\_base.h}{207}{IntegerLiteral}; this text says ``bound
literal'' throughout, but the two names refer to the same object.) It
\emph{negates cleanly}: over the integers, the negation of $x \ge 5$ is exactly
$x \le 4$. Moreover, a conjunction of bound literals describes an interval, so
the current domain $[\text{lb},\text{ub}]$ is just the two literals
$x \ge \text{lb}$ and $x \le \text{ub}$ held true at once. A variable is
\textbf{fixed} when its bounds meet, $\text{lb}(x) = \text{ub}(x)$: the interval
has collapsed to a single value, and this is the only sense in which $x$ ever
``has'' a value during search.

\textbf{Why this unifies the whole solver.} This representation lets a single
engine reason about a Boolean variable with domain $\{0,1\}$ and a wide integer
variable with domain $[0,10^6]$ through one common atom: the literal. A Boolean
assignment is already a literal, and an integer bound is a bound literal
$\bl{x \ge v}$ of the kind the previous paragraph introduced.
The two meet, but not because a Boolean is secretly an integer; it is the other
way around. When search needs to decide an integer bound, that bound is
\emph{reified}: turned into a Boolean literal (\cref{sec:encoding}) and handed to
the SAT core, which then drives the search on it: it can take that literal as a
decision, maintain its activity, and let learned clauses propagate it directly.
Deciding a Boolean and committing to an integer bound thus become the same kind
of step, recorded at \emph{one} shared decision level, which is precisely what
lets the SAT half and the CP half cooperate under \emph{one} conflict-analysis
mechanism rather than passing messages between two separate solvers.

\subsection{The trail}\label{sec:trail}
% Sources (VERIFIED):
%   ortools/sat/sat_base.h -- class Trail (331): the Boolean trail; AssignmentInfo
%     (277) stamps each entry with its decision level; EnqueueSearchDecision (425)
%     opens a new level for a decision; Reason() (525) returns a forced literal's
%     reason; Untrail() (595) truncates the trail on backjump.
%   ortools/sat/integer.h -- class IntegerTrail (523): the integer trail; Enqueue()
%     (745) pushes a bound change; integer_trail_ (1123) is the stack of bound
%     entries; Untrail() (547) unwinds it in lockstep with the Boolean trail.

A solver that guesses and backtracks needs a precise record of what it has
assumed and what followed. That record is the \textbf{trail}
(\srcla{sat\_base.h}{331}{Trail}): the ordered, append-only log of every literal
asserted on the current path from the root of the search tree to the node being
explored. Two kinds of entry land on it: \emph{decisions}, the literals the
solver guesses (\srcla{sat\_base.h}{425}{EnqueueSearchDecision} and
\cref{sec:complete}), and the literals propagation then forces from them.
Each entry is stamped with the \textbf{decision level} at which it was added
(\srcl{sat\_base.h}{277}), and each forced entry also carries a \emph{reason}
(\srcl{sat\_base.h}{525}), the facts that entailed it (\cref{sec:learning}).
Backtracking to level $k$ truncates the trail to its level-$k$ prefix
(\srcla{sat\_base.h}{595}{Untrail}), undoing every later assertion in one
step. The trail is thus both the solver's working memory and the substrate
conflict analysis walks backward when a guess fails.

CP-SAT reasons in two kinds of literal, so it keeps two coupled trails. Bound
changes accumulate on the \textbf{integer trail}
(\srcla{integer.h}{523}{IntegerTrail}, the stack \srcla{integer.h}{1123}{integer\_trail\_}), the stack of bound literals introduced above; Boolean
assignments accumulate on the \textbf{Boolean trail}. The two are not
independent: they share one decision-level counter, so they advance and unwind
together (the integer trail's \srcla{integer.h}{547}{Untrail} fires with
the Boolean one), and a bound the solver has reified into a Boolean is recorded
on both. An ordinary propagated bound, by contrast, lives only on the integer
trail. Keeping bound changes on a native integer trail, rather than as Booleans,
is what makes the encoding of \cref{sec:encoding} lazy: most bound movements
never touch the Boolean trail at all. Where this text says \emph{the} trail, it
means the single timeline the two structures form together.

\subsection{Propagation: the CP half}
% Sources (VERIFIED):
%   ortools/sat/integer.h -- PropagatorInterface (Propagate()/IncrementalPropagate());
%     GenericLiteralWatcher schedules registered propagators to a fixed point
%     (WatchLiteral/WatchLowerBound/WatchUpperBound, Register()).
%   Global constraints live in ortools/sat/ (e.g. all_different.*, no_overlap*/
%     scheduling*, circuit.*, cumulative.*, diffn.*).

With variables reduced to bounds, we can say precisely what a constraint
\emph{does} during search: it narrows those bounds. Each constraint comes with a
\textbf{propagator}, a piece of reasoning that, given the current bounds of the
variables it touches, tightens those bounds or reports a conflict when no value
remains.

Consider a simple example:
\[
  x \in [0, 10],\quad y \in [0, 10],\quad z \in [0, 10],
  \qquad \text{constraint: } x + y \le z .
\]
If, at some point, $z \le 4$ and $y \ge 2$, the linear propagator deduces
$x \le 2$ and pushes this new bound onto the trail. That deduction may wake
another propagator watching $x$, which may tighten another bound, and so on.
Propagators fire in cascade until the system reaches a \textbf{fixed point}: no
propagator can deduce anything new. Alternatively, one of them may detect a
\textbf{conflict}; for example, it may need $x \le 2$ while the trail already
contains $x \ge 3$, leaving the domain of $x$ empty.

Global constraints such as
\srcla{all\_different.h}{52}{AllDifferent}, \srcla{disjunctive.h}{49}{NoOverlap} (scheduling), \srcla{circuit.h}{46}{Circuit} (routing),
and \srcla{cumulative.h}{48}{Cumulative} (resources) carry specialized algorithms that prune far
more aggressively than a collection of generic inequalities could. Yet however
hard it prunes, \textbf{propagation alone never makes a choice}: everything it
concludes is a forced consequence of the current bounds. When it runs out of
necessary conclusions and the problem is not yet solved, search must guess.

\Cref{fig:sudoku-propagation} shows this guess-and-cascade on a Sudoku. One caveat
carries through the rest of this text: CP-SAT does \emph{not} carry hole-punched
domains like the pencil marks in that figure. It propagates only each variable's
lower and upper \emph{bounds} and hands every interior elimination to the SAT
engine as a Boolean literal. The bound-literal form of exactly this reasoning is
the implication graph in \Cref{fig:uipgraph}.

\subsection{Example: a propagator for $3x + 6y \le z$}\label{sec:onepropagator}
% Sources (VERIFIED):
%   ortools/sat/integer_expr.h:71-169 -- class LinearConstraintPropagator<use_int128>
%     : public PropagatorInterface, LazyReasonInterface; aliases IntegerSumLE /
%     IntegerSumLE128. Constraint stored canonically as sum_i coeff_i*var_i <= ub
%     (positive coeffs; a variable on the "other side" enters as NegationOf(var)).
%   integer_expr.cc:444-458 -- RegisterWith(): WatchLowerBound(vars_[i], id) for
%     every variable -> woken incrementally only when a watched LOWER bound moves.
%   integer_expr.cc:300-311 -- slack = ub - sum_i coeff_i*lb(var_i).
%   integer_expr.cc:334-368 -- Propagate(): for each i pushes new_ub_i = lb_i +
%     floor(slack/coeff_i) via EnqueueWithLazyReason(LowerOrEqual(var_i,new_ub),
%     0, propagation_slack, this).
%   integer_expr.cc:228-233 -- Explain() (LazyReasonInterface) returns ONLY the
%     raw reason: boolean_literals_at_false + vars_ + coeffs_ (ALL variables).
%     It does NOT read the slack, does NOT drop the explained var, does NOT relax.
%     (Args id/to_explain are unused in the body.)
%   integer.h:1145 LazyReasonEntry::Explain -- the TRAIL sets reason->slack =
%     propagation_slack (stored at EnqueueWithLazyReason time), then calls the
%     propagator's Explain().
%   integer.cc:1956-2022 ComputeLazyReasonIfNeeded -- the TRAIL drops the
%     propagated var (to_ignore) and, if slack>0, calls RelaxLinearReason(
%     reason->slack, coeffs, indices) at integer.cc:2021. THIS is where the
%     stored slack is consumed for an ordinary push.
%   integer_expr.cc:110-121 -- FillIntegerReason(): collects the bounds used.
%   integer_expr.cc:299-327 -- CONFLICT path only (slack<0): the propagator
%     itself calls RelaxLinearReason(-slack-1,...) at :314 + ReportConflict() at
%     :325 (or EnqueueLiteral when a single enforcement literal is unassigned).

In the CP-SAT source a propagator is wrapped in templates, watch lists, and
lazy-reason plumbing, but the object underneath is small. As a worked example, consider the linear propagator for $3x + 6y \le z$ over
$x, y \in [0,3]$ and $z \in [0,30]$, following what CP-SAT's
\srcla{integer\_expr.h}{72}{LinearConstraintPropagator} actually does. CP-SAT stores
every linear constraint in \emph{canonical form}, all terms on one side of a
constant, so throughout we work with it as
\[
  3x + 6y - z \le 0 ,
\]
where the term $-z$ is the variable $\srcla{integer\_base.h}{175}{NegationOf}(z)$.
The role comes with two duties. In a classical finite-domain
solver a propagator has only the first: narrow domains, then say no more. A
CP-SAT propagator must also \emph{justify} each deduction on demand, naming the
bounds that forced it, so that the learning half (\cref{sec:learning}) can
distill a failure into a permanent rule. That second duty, explaining a
propagation rather than merely performing it, is what \emph{lazy clause
generation} adds to ordinary propagation; it is the one a classical propagator
does not carry.

\textbf{We are called only when a watched bound changes.} Running the propagator on
every change in the search would be wasteful; we want it to fire only when something
has happened that could let us deduce more. So on creation we register with the
\srcla{integer.h}{1324}{GenericLiteralWatcher}, asking to be called whenever the
\emph{lower} bound of $x$, $y$, or $-z$ changes (\srcla{integer\_expr.cc}{449}{WatchLowerBound} on each); we are never run on a step that touched none of the
three. For a ``$\le$'' constraint only the lower bounds matter: they are what can
push the left-hand side up toward the bound. Because $z$ enters as $-z$, the update
``$z$'s upper bound fell to $8$'' reaches us as ``the lower bound of $-z$ rose to
$-8$,'' so a single watch direction already covers both sides of every variable,
directly reflecting the bounds-as-everything data model of \cref{sec:bounds}.

\textbf{If the constraint is violated, we report a conflict.} Whenever a watched
bound moves, we recompute our \emph{slack}, the room left before the constraint
binds (\srcl{integer\_expr.cc}{310}):
\[
  \text{slack} \;=\; \underbrace{0}_{\text{rhs}}
  \;-\; \big(3\,\text{lb}(x) + 6\,\text{lb}(y) + 1\cdot\text{lb}(-z)\big) .
\]
If it is negative, the current bounds already rule the constraint out and there is
nothing to deduce: we report a conflict. Suppose decisions have forced $x \ge 2$
and $y \ge 1$ while $z \le 8$, so $\text{lb}(x) = 2$, $\text{lb}(y) = 1$, and
$\text{lb}(-z) = -8$; then
\[
  \text{slack} \;=\; 0 - (3\cdot 2 + 6\cdot 1 + 1\cdot(-8)) \;=\; -4 \;<\; 0
\]
(\srcl{integer\_expr.cc}{312}). We do not backtrack or learn; that is the SAT
half's job. We simply fill our reason $\{x \ge 2,\ y \ge 1,\ z \le 8\}$ and call
\srcla{integer\_expr.cc}{325}{ReportConflict}, declaring that these bounds
are jointly impossible. The conflict analysis of \cref{sec:learning}
takes it from there.

\textbf{What we read, and what we push.} When the slack is non-negative there is
room to spare, and instead of failing we tighten bounds. Take the lighter state
where $x$ and $y$ still sit at their lower bound $0$ while $z \le 8$, so the slack
is $0 - (3\cdot 0 + 6\cdot 0 + 1\cdot(-8)) = 8$.
This is not a conflict, but we could already prove that $x \le 2$ and $y \le 1$ must hold, since any higher would push the left-hand side above $8$.
Reading only the \emph{lower} bounds of the other
variables, we push an \emph{upper} bound on each variable
(\srcl{integer\_expr.cc}{334}):
\[
  \text{new\_ub}_i =
  \text{lb}_i + \left\lfloor \frac{\text{slack}}{\text{coeff}_i} \right\rfloor .
\]
For this example, the resulting bounds are
\[
  x \le 0 + \lfloor 8/3 \rfloor = 2, \qquad\qquad
  y \le 0 + \lfloor 8/6 \rfloor = 1 .
\]
Both tighten the domain: $x$ shrinks from $[0,3]$ to $[0,2]$, and $y$ shrinks
from $[0,3]$ to $[0,1]$. That is the whole job: read the other variables' bounds
and push tighter bounds of our own. There is no guessing and no search, only
forced consequences.
% Source (VERIFIED): template LinearConstraintPropagator<use_int128> and the alias
%   IntegerSumLE128 (integer_expr.h:71,168-169); int128 activity/slack at
%   integer_expr.cc:258-310 ("absl::int128 lb_128 = 0;", "slack128 = int128(ub) - lb_128").
\begin{platypusinfo}[htb]
When a constraint's coefficients and bounds are large enough that the running
activity could overflow $64$ bits, CP-SAT instantiates the linear-sum propagator
in a $128$-bit variant (\srcla{integer\_expr.h}{169}{IntegerSumLE128}): the
coefficients remain \code{int64}, but activity and slack accumulate in
\code{int128}. The bound it pushes is therefore the exact integer bound, not a
wrapped-around value: the solver keeps its arithmetic exact all the way down to a
single linear row.
\end{platypusinfo}

\textbf{Why do we not directly add these insights as clauses?}
In principle we could add a clause for the implication
\[
  (\bl{x \ge 0} \;\wedge\; \bl{z \le 8}) \;\Rightarrow\; \bl{y \le 1}
\]
outright, but that is both costly and unnecessary. For now it is enough to tell the
search that $y \le 1$, which updates the bounds and may wake other propagators; we
owe an explanation only if $\bl{y \le 1}$ later turns out to be part of a conflict, when
conflict analysis needs the deduction to trace the root cause and learn a clause.
And the more precise that reason, the more effective the learning: here it is
trivial, but for a constraint over many variables it can pay to do real work at
explain time and pinpoint exactly the subset of bounds that forced the deduction.
CP-SAT's \srcla{cp\_constraints.cc}{123}{GreaterThanAtLeastOneOf} propagator does
exactly this: its \code{Explain} drops the bounds already implied at the root and
returns that trimmed subset rather than the whole row. Since most deductions never
enter a conflict at all (look at how few bounds the conflict in
\cref{fig:sudoku-conflict} actually rests on), computing and storing a precise
reason for every one of them would be wasted effort.

\textbf{So how is the reason attached?} Rather than record the justification, we
call \srcla{integer\_expr.cc}{360}{EnqueueWithLazyReason} and hand the trail
only a pointer back to ourselves. The trail enqueues $y \le 1$ and wakes whatever
that unblocks, but stores no explanation alongside it. Only if this deduction is
later dragged into a conflict does the engine call our \srcla{integer\_expr.cc}{228}{Explain} method; only then do we name the bounds that forced the push,
$x \ge 0$ and $z \le 8$, handing back the implication above for conflict analysis
(\cref{sec:learning}) to resolve and learn from. The reason is available the
whole time, but built only for the deductions that turn out to need it. This is lazy
clause generation at the scale of a single propagator.

At its core, the entire object consists of two short
procedures: one that pushes bounds and one that explains them on request. The
trail owns the events, the propagator only responds, and the explanation is a
deferred call-back that the trail makes only if a pushed bound is later dragged
into a conflict.

\begin{algorithm}[htbp]
\caption{Linear propagator for $\sum_i c_i x_i \le u$ in canonical form (all $c_i > 0$)}
\label{alg:linprop}
\begin{algorithmic}[1]
\Procedure{Propagate}{} \Comment{woken when a watched lower bound rises}
  \State $\text{slack} \gets u - \sum_i c_i \cdot \text{lb}(x_i)$ \Comment{room before the constraint binds}
  \If{$\text{slack} < 0$}
    \State \Return \Call{ReportConflict}{$\{\, \bl{x_i \ge \text{lb}(x_i)} \,\}_i$} \Comment{bounds jointly impossible}
  \EndIf
  \ForAll{$i$}
    \State $\text{newub} \gets \text{lb}(x_i) + \lfloor \text{slack}/c_i \rfloor$
    \If{$\text{newub} < \text{ub}(x_i)$}
      \State \Call{EnqueueWithLazyReason}{$\bl{x_i \le \text{newub}}$, \ \textbf{self}}
    \EndIf
  \EndFor
\EndProcedure
\Statex
\Procedure{Explain}{$\bl{x_j \le v}$} \Comment{called \emph{only} if this pushed bound enters a conflict}
  \State \Return $\{\, \bl{x_i \ge \text{lb}(x_i)} \,\}_i$ \Comment{the raw reason: \emph{all} variables, unrelaxed}
\EndProcedure
\end{algorithmic}
\end{algorithm}

\subsection{Lifting the reason}\label{sec:relaxation}

We kept the propagator deliberately simple: its \code{Explain} hands back the raw
row, the full set of bounds that were in force, without hunting for the smallest or
most general subset. CP-SAT makes up for that afterward. Conflict analysis
(\cref{sec:learning}) turns each reason into a permanent clause, and a
\emph{weaker} reason, one that holds in more states, becomes a clause that rules out
more of them. So before learning, CP-SAT \emph{lifts} the reason
(\srcla{integer.cc}{1083}{RelaxLinearReason}): it spends the slack the deduction
left unused to loosen the bounds, recovering a more general clause than the
propagator literally reported.

Why is this sound? Because a push rarely uses all the room it had. Take the push of
$y \le 1$, made when the slack was $8$; reading off $\lfloor 8/6 \rfloor = 1$ gave
$y \le 1$, but the leftover
\[
  \ell \;=\; \big(\lfloor 8/6 \rfloor + 1\big)\cdot 6 \;-\; 8 \;-\; 1
        \;=\; 2\cdot 6 - 8 - 1 \;=\; 3
\]
measures how much room was to spare. Forcing $y \le 1$ never needed the full bound
$z \le 8$: even $z \le 11$ leaves $y \ge 2$ impossible once $x \ge 0$ is known,
since $3\cdot 0 + 6\cdot 2 = 12 > 11$. So CP-SAT lifts $z \le 8$ to $z \le 11$,
weakening the bound by exactly $\ell = 3$, and conflict analysis resolves against
the broader implication
\[
  (x \ge 0 \;\wedge\; z \le 11) \;\Rightarrow\; y \le 1
\]
instead of the tighter $z \le 8$ form, learning a clause that can fire in more
future states. The same trick broadens reported conflicts: there the slack is
already negative, and its magnitude is the budget for loosening the bounds the
violation did not strictly need.

All the propagator had to store to make this possible was a single integer, the
leftover slack, carried alongside its claim check rather than a copied reason. That
one number is what lets CP-SAT recover not merely \emph{a} reason but a deliberately
general one.

\subsection{Expansion: which constraints keep a propagator}\label{sec:expansion}
% Sources (VERIFIED):
%   ortools/sat/cp_model_expand.cc -- ExpandElement (1159), ExpandPositiveTable
%     (2288)/ExpandNegativeTable (1673), ExpandAutomaton (1318), ExpandInverse (679),
%     MaybeExpandAllDiff (2713, heuristic: small / near-permutation / fully encoded),
%     ExpandReservoir (406, default-ON via expand_reservoir_constraints=true),
%     ExpandLinMax (808, size-gated; max_lin_max_size_for_expansion default 0 = kept). Default behavior
%     KEEPS native propagators; expansion is heuristic / parameter-gated.
%   ortools/sat/cp_model_loader.cc -- kept propagators are loaded here:
%     LoadBoolOr/BoolAnd/BoolXor (clauses), LoadLinearConstraint (1213, supports
%     enforcement literals = half-reification), LoadAllDiffConstraint -> AllDifferentOnBounds
%     (1485), LoadCircuit/Routes (1664/1679), LoadNoOverlap/Cumulative/NoOverlap2d
%     (1610/1629/1616), LoadIntProd/IntDiv/FixedModulo (1506/1553/1576).
%   VISIBLE LINKS for kept constraints point past the one-line loader to the
%     algorithm header where the propagator/factory + its doc comment live
%     (re-verify these line numbers when re-pinning \orsha):
%       AllDifferent (kept "on bounds") -> all_different.h:52 (AllDifferentOnBounds)
%       NoOverlap    -> disjunctive.h:49 (AddDisjunctive)
%       NoOverlap2d  -> diffn.h:89 (AddNonOverlappingRectangles)
%       Cumulative   -> cumulative.h:48 (Cumulative)
%       Circuit      -> circuit.h:46 (class CircuitPropagator)
%       MultipleCircuit -> circuit.h:259 (LoadSubcircuitConstraint, wraps CircuitPropagator)
%     The heuristically-expanded AllDifferent mention still links MaybeExpandAllDiff
%     (cp_model_expand.cc:2713); the Expanded-away names keep their Expand* targets.

\cref{sec:onepropagator} said that every constraint type has a
propagator. That is only half the story. A modeling language offers dozens of
high-level constraints, and CP-SAT does \emph{not} keep a bespoke propagator for
each of them. During presolve, many constraints are \textbf{expanded}: rewritten
into the small set of primitives that the engine reasons about natively, such
as clauses, linear constraints, at-most-one constraints, and value literals.
Only a curated subset is \textbf{kept} as dedicated global propagators.

\textbf{Expanded away} (\srcf{cp\_model\_expand.cc}):
\srcla{cp\_model\_expand.cc}{1159}{Element},
\srcla{cp\_model\_expand.cc}{2288}{Table} (positive and negative),
\srcla{cp\_model\_expand.cc}{1318}{Automaton}/regular, and
\srcla{cp\_model\_expand.cc}{679}{Inverse}. CP-SAT may also expand
\srcla{cp\_model\_expand.cc}{2713}{AllDifferent} heuristically when its domain
is small or close to a permutation; in that case, it becomes value literals plus
at-most-one constraints. Two more sit on a parameter-controlled line:
\srcla{cp\_model\_expand.cc}{406}{Reservoir} is expanded \emph{by default} (a
parameter can instead keep it native), while
\srcla{cp\_model\_expand.cc}{808}{LinMax} constraints keep their propagator
\emph{by default}, with a size threshold that opts only the \emph{small} ones into
expansion.
% LinMax gating (cp_model_expand.cc:2875): expand iff exprs().size() <=
%   max_lin_max_size_for_expansion, default 0 (sat_parameters.proto:578) -> nothing
%   expanded unless the threshold is raised, and then it is the SMALL constraints that
%   get expanded, not the large ones. Prose previously said "large ... unless a size
%   threshold opts them into expansion", which points the mechanism at the wrong end.
The common thread is not
that a stronger native propagator could never exist; it is that the decomposition
is often cheap, exposes useful Boolean or value structure to the rest of the
engine, or is competitive enough that a dedicated algorithm would not reliably
pay for itself.

\textbf{Kept as native propagators} (\srcf{cp\_model\_loader.cc}): Boolean
constraints, loaded as clauses; \textbf{linear} constraints, crucially
\emph{with enforcement literals}, that is, half-reified ``indicator'' constraints
in which a literal $\ell$ enforces a linear constraint; integer \textbf{product}
(variable~$\times$~variable), \textbf{division} (the divisor may be a variable),
and \textbf{modulo} constraints (the modulus must be a fixed constant; a variable
modulus is decomposed into division and product during presolve);
% int_prod/int_div/int_mod loaded natively in cp_model_loader.cc (VERIFIED):
%   LoadIntProd (1506) -> ProductConstraint (integer_expr.h:793); fully general var x var.
%   LoadIntDiv  (1553): denom fixed -> FixedDivisionConstraint (integer_expr.h:836),
%     ELSE DivisionConstraint (integer_expr.h:817). A variable divisor IS native --
%     so "fixed divisors" was wrong for division (and meaningless for product).
%   LoadIntMod  (1576): CHECK(IsFixed(mod)) -> FixedModuloConstraint (integer_expr.h:851);
%     there is no variable-modulus propagator. ExpandIntMod (cp_model_expand.cc:575)
%     rewrites a non-fixed modulus into int_div + int_prod + linear during presolve.
\srcla{all\_different.h}{52}{AllDifferent} on bounds when it is \emph{not} a
small permutation; \srcla{circuit.h}{46}{Circuit} and
\srcla{circuit.h}{259}{MultipleCircuit} for routing; and the scheduling trio
\srcla{disjunctive.h}{49}{NoOverlap}, \srcla{cumulative.h}{48}{Cumulative}, and
\srcla{diffn.h}{89}{NoOverlap2d}. These are the constraints whose specialized
algorithms (edge-finding, energetic reasoning, subtour elimination, and
Hall-interval pruning) can prune far more than any clause decomposition could.
This is where the ``CP'' in CP-SAT does its real work.

This point is easy to oversimplify, so one nuance is worth stating plainly: the
released code \emph{defaults to keeping} native propagators, and expansion is
gated by size heuristics and parameters rather than applied wholesale. The
reason ties directly back to the cost argument of this part. A dedicated global
propagator is expensive, so CP-SAT keeps it only where its extra pruning tends
to pay for itself; everything else is handed to the cheap, always-on clause and
linear machinery.
\subsection{Search: guessing when propagation stalls}\label{sec:complete}
% Sources (VERIFIED):
%   ortools/sat/integer_search.h -- BooleanOrIntegerLiteral (a decision is a
%     Boolean OR an integer bound literal).
%   ortools/sat/sat_solver.cc -- EnqueueDecisionAndBackjumpOnConflict (binary
%     decisions, backjump on conflict). Heuristic detail moved to sec:search.

Whether a constraint reasons through a native propagator or an expanded
decomposition, propagation eventually runs dry: it reaches a fixed point with
nothing left to deduce. If every variable happens to be fixed, that fixed point
is a solution. Usually it is not. With no conflict, the fixed point may still
leave $3 \le x \le 7$ open: if every constraint is satisfiable across that whole
range, no propagator has any reason to tighten it. That is a consistent
\emph{state}, but not a \emph{solution}. A solution must assign $x$ one concrete
value, and propagation alone, which only ever draws forced conclusions, cannot
supply it. Closing the gap falls to \emph{search}.

The solver's only move at a stalled fixed point is to \emph{guess}. The guess is
a \textbf{decision}: it asserts a single literal, for instance trying
$\bl{x \le 7}$ true, with its negation $\bl{x \ge 8}$ held in reserve as the other branch. As \cref{sec:bounds}
noted, an integer bound is itself a reified literal, so guessing
$\bl{x \le 7}$ and flipping a native Boolean are the same step,
both recorded at a new \textbf{decision level}.

A decision sets off a fresh cascade of propagation, and the guess is judged by
what that cascade returns: it may pin more variables and let the solver descend
with another guess; it may fix everything with no conflict, giving a solution; or
it may reach a conflict, proving the guess wrong and forcing the solver to
\emph{backtrack} and learn (\cref{sec:learning}). Decisions stack on the
\textbf{trail}, each tagged with its decision level, and backtracking pops them
off and relaxes the bounds they imposed. Search is then just this loop: guess a
literal, let the resulting cascade say whether it was right, wrong, or simply not
yet decisive, and guess again until every variable is pinned. \emph{How} the
solver picks each literal, and what it does once the Booleans are all assigned
but an integer still spans a range, is the search machinery assembled in
\cref{sec:search}.
